\section{System design}\label{sec:design}

This section specifies the processing chain in the order a recording passes
through it, and then the four classifiers placed at its end.
Section~\ref{sec:d-overview} gives the shape of the pipeline and fixes
notation. Section~\ref{sec:d-cond} covers audio conditioning.
Section~\ref{sec:d-wav} develops the wavelet denoising stage.
Section~\ref{sec:d-spec} covers the two spectrogram representations.
Section~\ref{sec:d-arch} describes the four architectures, and
Section~\ref{sec:d-why} states why this particular set was chosen.

\subsection{Overview and notation}\label{sec:d-overview}

Let $x[n]$ denote a recording sampled at rate $f_s$, and let
$y \in \{\text{cough}, \text{noncough}\}$ denote the label to be predicted.
The pipeline applies four operations in sequence. The waveform is brought to a
common sampling rate and channel count; it is optionally denoised in a wavelet
basis; it is converted to a time frequency image; and the image is classified.
Figure~\ref{fig:pipeline} shows the arrangement, including the branch point at
which the two spectrogram representations diverge and the optional stage whose
contribution is the question this paper poses.

The optional stage is what makes the design a factorial one rather than a
single pipeline. Writing $\mathcal{W}$ for the denoising operator,
$\mathcal{S}$ for the spectrogram operator and $g_\theta$ for a classifier
with parameters $\theta$, the eight configurations evaluated in
Section~\ref{sec:results} are
\begin{equation}
\hat{y} = g_\theta\big(\mathcal{S}(\mathcal{W}^{\,\delta}(x))\big),
\qquad \delta \in \{0, 1\},
\label{eq:config}
\end{equation}
with $\delta = 1$ indicating that denoising is applied and $\delta = 0$ that
it is bypassed, and with $\mathcal{S}$ and $g_\theta$ ranging over the
representations and architectures described below. Holding everything outside
Equation~\eqref{eq:config} fixed is what allows the effect of $\delta$ to be
attributed to denoising rather than to an incidental difference in the
experimental conditions.

\begin{figure}[t]
\centering
\begin{tikzpicture}[
  node distance=6mm,
  box/.style={draw, rounded corners=2pt, align=center, inner sep=4pt,
              minimum height=9mm, font=\footnotesize},
  opt/.style={box, dashed},
  ar/.style={-{Latex[length=2mm]}, thick}
]
\node[box] (wav) {Recording\\(WAV)};
\node[box, right=of wav] (rs) {Resample 16~kHz\\mono};
\node[opt, right=of rs] (wd) {Wavelet denoise\\sym4, $J=6$};
\node[box, right=of wd, yshift=8mm] (stft) {Linear axis\\spectrogram};
\node[box, right=of wd, yshift=-8mm] (mel) {Perceptual axis\\spectrogram};
\node[box, right=of stft] (cnn) {Compact network\\(from scratch)};
\node[box, right=of mel] (tl) {ResNet18\\MobileNetV2\\DenseNet201};
\node[box, right=of cnn, yshift=-8mm] (out) {Cough /\\noncough};

\draw[ar] (wav) -- (rs);
\draw[ar] (rs) -- (wd);
\draw[ar] (wd.east) -- ++(3mm,0) |- (stft.west);
\draw[ar] (wd.east) -- ++(3mm,0) |- (mel.west);
\draw[ar] (stft) -- (cnn);
\draw[ar] (mel) -- (tl);
\draw[ar] (cnn.east) -- ++(3mm,0) |- (out.west);
\draw[ar] (tl.east) -- ++(3mm,0) |- (out.west);
\draw[ar] (rs.south) -- ++(0,-20mm) -| (mel.south);
\end{tikzpicture}
\caption{Processing chain. The denoising stage is drawn dashed because it is
present in half of the configurations and absent in the other half, which is
the comparison the study is built around. The lower path from the resampling
stage carries the undenoised signal directly to spectrogram extraction.
Beyond the stage order, the figure shows that the two spectrogram
representations are formed from the same conditioned waveform, so any
difference in the results attaches to the frequency axis rather than to
anything earlier.}
\label{fig:pipeline}
\end{figure}

\subsection{Audio conditioning}\label{sec:d-cond}

Recordings arrive at heterogeneous sampling rates and channel counts, and both
are normalised before anything else is done. Files recorded in stereo are
reduced to a single channel by averaging, which preserves content common to
both channels and attenuates anything uncorrelated between them. All material
is then resampled to $f_s = 16$~kHz.

The choice of rate is a compromise between fidelity and cost. A rate of
16~kHz places the Nyquist limit at 8~kHz, which retains the spectral region in
which cough energy is concentrated while discarding content above it, and it
reduces the sample count of a typical recording by roughly two thirds relative
to the 44.1~kHz material that consumer devices commonly produce. The explosive
onset of a cough does deposit some energy above 8~kHz, so the choice is not
free of cost; the justification is that the discriminative information is
dominated by the lower range and that the reduction in sample count propagates
through every subsequent stage, including the wavelet decomposition, whose
cost is linear in signal length.

\subsection{Wavelet denoising}\label{sec:d-wav}

Denoising is performed in a wavelet basis rather than by filtering in the
frequency domain, for a reason specific to the signal. A cough is a transient,
so its energy is localised in time and spread across frequency, which is the
opposite of the profile that a fixed frequency domain filter handles well. A
wavelet basis represents the signal at several resolutions at once and so can
distinguish a short broadband event from sustained interference occupying the
same band, which a filter cannot.

The decomposition expresses the signal as a coarse approximation at level $J$
together with the detail removed at every level above it,
\begin{equation}
x[n] = A_J[n] + \sum_{j=1}^{J} D_j[n],
\label{eq:mra}
\end{equation}
where $A_J[n]$ holds the low frequency content remaining after $J$ successive
halvings of the band and $D_j[n]$ holds the detail at level $j$. Each
level separates the current approximation into a new approximation and a
detail band, so $D_1$ carries the highest frequency octave and $D_J$ the
lowest of those removed. We set $J = 6$ and use the sym4 wavelet, a member of
the symlet family chosen because it is nearly symmetric, which limits the
phase distortion introduced into the reconstructed transient, and because its
support is short enough to remain localised against events lasting a few tens
of milliseconds.

Denoising proceeds by attenuating detail coefficients whose magnitude is
consistent with noise rather than with signal. The classical treatment fixes
a threshold per level and shrinks each coefficient toward zero by that
amount, discarding those that do not survive the
subtraction~\cite{donoho1994, donoho1995}. The procedure used here is the
empirical Bayesian alternative, which places a mixture prior on the detail
coefficients at each level,
\begin{equation}
d_{j,k} \sim (1 - w_j)\,\delta_0 + w_j\,\gamma,
\label{eq:eb}
\end{equation}
where $\delta_0$ is a point mass at zero, $\gamma$ is a heavy tailed Cauchy
density and the mixing weight $w_j$ is estimated from the coefficients at
level $j$ by marginal maximum likelihood. Each coefficient is then replaced
by the median of its posterior distribution~\cite{johnstone2005}. Because
$w_j$ is estimated separately at every level, the amount of shrinkage adapts
to how much of the energy at that scale is attributable to noise instead of
being set by one global rule, and the heavy tail of the prior leaves large
coefficients close to unshrunk. The signal is reconstructed from the retained
approximation and the shrunken details.

Figure~\ref{fig:denoise} shows the operation on one cough and one noncough
recording, each at three stages. Reading across the top row establishes that
the transient structure of the cough survives the operation, which is the
property the stage must have if it is to be useful; reading down the third
column compares what was discarded in the two cases, which is where the two
classes differ most visibly.

\begin{figure}[t]
\centering
\begin{tabular}{@{}ccc@{}}
\includegraphics[width=0.30\textwidth]{denoise_a_cough_raw.png} &
\includegraphics[width=0.30\textwidth]{denoise_b_cough_denoised.png} &
\includegraphics[width=0.30\textwidth]{denoise_c_cough_removed.png} \\[2pt]
{\footnotesize (a) cough, original} &
{\footnotesize (b) cough, denoised} &
{\footnotesize (c) cough, removed} \\[6pt]
\includegraphics[width=0.30\textwidth]{denoise_d_noncough_raw.png} &
\includegraphics[width=0.30\textwidth]{denoise_e_noncough_denoised.png} &
\includegraphics[width=0.30\textwidth]{denoise_f_noncough_removed.png} \\[2pt]
{\footnotesize (d) noncough, original} &
{\footnotesize (e) noncough, denoised} &
{\footnotesize (f) noncough, removed} \\
\end{tabular}
\caption{Effect of the sym4 decomposition at level six on a representative
cough recording, in the top row, and a representative noncough recording, in
the bottom row. Each panel is a time frequency representation, with time on
the horizontal axis and frequency on the vertical. Panels (a) and (d) show
the conditioned input and (b) and (e) the reconstruction after shrinkage. The
dominant transient structure of the cough is preserved from (a) to (b), which
is the property the operation has to have if it is not to suppress the energy
the classifier depends on. Panels (c) and (f) are diagnostic rather than
quantitative: each is the absolute difference between the two images above it,
rescaled to the full display range and amplified to make faint structure
visible. They indicate where in the plane the reconstruction departs from the
input, and that the departure avoids the region carrying the cough transient.
Because each difference panel is rescaled independently of the other, the
intensities in (c) and (f) are not on a common scale and the two panels should
not be compared against each other.}
\label{fig:denoise}
\end{figure}

The level $J = 6$ determines how much of the signal is exposed to shrinkage.
At 16~kHz the sixth level corresponds to an approximation band below roughly
125~Hz, so the decomposition reaches well beneath the range in which cough
energy is concentrated and gives the thresholding access to the bands where
both the transient and the interference reside. A shallower decomposition
would leave low frequency interference untouched in the approximation; a
deeper one would continue subdividing a band that by then contains little but
the recording offset.

\subsection{Spectrogram representations}\label{sec:d-spec}

Both representations begin from the same conditioned waveform and differ only
in how the frequency axis is spaced, which is what allows their comparison to
isolate the effect of that spacing.

\subsubsection{Linear frequency axis}\label{sec:d-stft}

The short time Fourier transform is computed as
\begin{equation}
X(m, k) = \sum_{n} x[n]\, w[n - m]\, e^{-j 2 \pi k n / N},
\label{eq:stft}
\end{equation}
where $m$ indexes the analysis frame, $k$ the frequency bin, $w[\cdot]$ is a
Hamming window and $N$ is the transform length. The magnitude $|X(m,k)|$ is
retained and the phase discarded, since the classifier operates on an image.

The representation preserves the full time frequency structure of the signal
but spaces its bins uniformly, which distributes resolution evenly across the
spectrum. For a cough this allocation is inefficient. The region above 4~kHz
occupies half of the available bins while carrying comparatively little of
what separates a cough from a door closing, and the region below 500~Hz, where
the voiced termination and much of the envelope structure reside, is
represented by a small number of bins. This representation is used for the
baseline network described in Section~\ref{sec:d-cnn}.

\subsubsection{Perceptual frequency axis}\label{sec:d-mel}

The perceptual representation warps the frequency axis according to
\begin{equation}
f_{\mathrm{mel}} = 2595 \log_{10}\!\left(1 + \frac{f}{700}\right),
\label{eq:mel}
\end{equation}
in which $f$ is the physical frequency in hertz and $f_{\mathrm{mel}}$ the
warped value~\cite{davis1980}. The mapping is approximately linear below
about 700~Hz and approximately logarithmic above it, so the low frequency
region is expanded relative to a linear axis and the high frequency region
compressed. A bank of triangular filters spaced uniformly on this warped axis
is applied to the magnitude spectrum, and the filter outputs are converted to
a logarithmic amplitude scale.

Figure~\ref{fig:mel} shows the mapping and the bank it produces. Panel (a)
plots Equation~\eqref{eq:mel} against a linear reference, and the departure
between the two curves is the redistribution of resolution the mapping
performs. Panel (b) shows the consequence for the filters themselves, which
is the quantity that matters for what the network receives.

\begin{figure}[t]
\centering
\includegraphics[width=\textwidth]{mel_mapping.png}
\caption{The perceptual frequency mapping of Equation~\eqref{eq:mel} in panel
(a), plotted against a linear reference over the band retained after
resampling, and in panel (b) a reduced bank of sixteen triangular filters
spaced uniformly on the warped axis and drawn against hertz. Beyond restating
the mapping, the figure shows what it costs and buys: filters near the bottom
of the band span on the order of 240~Hz while those near the top span on the
order of 2200~Hz, so fine structure above 4~kHz is deliberately pooled away
in exchange for resolution below 1~kHz. The bank used in the experiments has
128 filters rather than the sixteen drawn here, which are reduced in number
only so the individual triangles remain distinguishable.}
\label{fig:mel}
\end{figure}

The triangular weightings overlap, each beginning at the centre of the one
before it and ending at the centre of the one after, so that every frequency
contributes to exactly two filters and the resulting representation is smooth
in the filter index. Because the filters are uniformly spaced on the warped
axis, they become progressively wider in hertz as frequency rises, which is
what performs the compression: a filter near 4~kHz pools energy over a band
several times wider than one near 300~Hz, and the fine structure within the
wider band is deliberately discarded. The logarithmic amplitude scaling
applied afterwards compresses the dynamic range so that the low energy
portions of a recording, including the decaying tail of a cough, remain
visible alongside the high energy onset rather than being driven to the floor
of the image.

We use a Hamming window of 512 samples with 256 samples of overlap, a
transform length of 512 and 128 filters, and resize the resulting image to
$224 \times 224$ pixels to match the input expected by the transferred
architectures. At 16~kHz a 512 sample window spans 32~ms and advances by
16~ms, which is short enough to resolve the onset of a cough as a distinct
event and long enough to give usable frequency resolution. The resize is
performed after the logarithmic scaling rather than before, so that
interpolation acts on a quantity whose dynamic range has already been
compressed.

\subsection{Classification architectures}\label{sec:d-arch}

Four architectures are evaluated. One is trained from random initialisation
and three are transferred from a large image corpus with the classification
head replaced by a two unit layer. Table~\ref{tab:arch} summarises the
properties that bear on the comparison.

\begin{table}[t]
\centering
\footnotesize
\caption{The four architectures and the properties relevant to the
comparison. Depth is counted as weighted layers along the principal path. The
final column records the mechanism each architecture uses to propagate
information across depth, which is the property this study relates to the
benefit obtained from denoising.}
\label{tab:arch}
\begin{tabular}{@{}llll@{}}
\toprule
Architecture & Input & Initialisation & Propagation mechanism \\
\midrule
Compact network & Linear axis & Random & Plain stacked convolutions \\
ResNet18 & Perceptual axis & Transferred & Additive identity shortcuts \\
MobileNetV2 & Perceptual axis & Transferred & Inverted residual bottlenecks \\
DenseNet201 & Perceptual axis & Transferred & Concatenative dense reuse \\
\bottomrule
\end{tabular}
\end{table}

\subsubsection{Compact network trained from scratch}\label{sec:d-cnn}

The baseline comprises three convolutional layers, each followed by a pooling
operation that halves the spatial dimensions, and two fully connected layers
ending in a two unit output. It is trained from random initialisation on the
linear axis spectrograms and establishes what is achievable without
transferred weights. Its role is to separate two effects that are otherwise
confounded: the gain from the perceptual frequency axis and the gain from
transferred initialisation. Because this model differs from the others in both
respects, it is treated throughout as a reference point rather than as a
competitor, and the comparisons that carry the argument of the paper are made
among the three transferred models.

\subsubsection{Residual network}\label{sec:d-resnet}

ResNet18 introduces additive shortcut connections that carry the input of a
block forward to its output, so that each block learns a residual correction
rather than a complete transformation~\cite{he2016}. The arrangement allows
gradients to reach early layers without attenuation through the intervening
nonlinearities, which is what permits depth to be increased without the
degradation that plagued plain stacks. At eighteen weighted layers it is the
shallowest of the three transferred models, and its capacity is moderate by
current standards, which is the property that makes it informative here.

\subsubsection{Inverted residual network}\label{sec:d-mobilenet}

MobileNetV2 is designed for deployment under tight memory and power budgets.
It replaces standard convolutions with depthwise separable ones, factorising a
convolution into a spatial filtering step applied per channel and a pointwise
combination across channels, which reduces the parameter count and the
arithmetic by roughly an order of magnitude for a given receptive field. Its
blocks expand the channel dimension internally and project back to a narrow
linear bottleneck, with the shortcut connecting the bottlenecks rather than the
expanded representations~\cite{sandler2018}. Including it tests whether a
detector of practical size on portable hardware can approach the accuracy of a
larger model, which matters because a cough detector is most useful where it
can run continuously.

\subsubsection{Densely connected network}\label{sec:d-densenet}

DenseNet201 connects every layer within a block to all preceding layers by
concatenation rather than addition, so a layer has direct access to every
feature map produced before it~\cite{huang2017}. The arrangement strengthens
gradient flow and reduces redundancy, since a feature computed once need not
be recomputed downstream. It is by a wide margin the deepest of the models
considered. Its inclusion is deliberate rather than incidental: if the benefit
of denoising diminishes as a network becomes better able to form
representations that ignore interference, then the architecture with the most
extensive feature reuse is where that diminution should be visible, and
Section~\ref{sec:res-denoise} reports what is observed.

\subsection{Rationale for the architecture set}\label{sec:d-why}

The four models were not selected for coverage of the literature but to span
the mechanisms by which a network can achieve robustness to interference. The
compact network has no such mechanism beyond pooling. ResNet18 propagates
information additively through a moderate depth. MobileNetV2 is explicitly
constrained in capacity, so whatever robustness it attains must come from
structure rather than from parameters. DenseNet201 maximises feature reuse.
If denoising substitutes for architectural robustness, as
Section~\ref{sec:rel-denoise} argues it should, then its benefit ought to
decline along this ordering, and the set is arranged so that the prediction
can be tested.

\subsection{Computational considerations}\label{sec:d-cost}

The pipeline is intended for a setting in which audio arrives continuously, so
the cost of each stage is worth stating, both because it bounds what can run
on portable hardware and because it determines whether the optional denoising
stage is worth its place.

Conditioning is linear in the number of samples. For a recording of duration
$T$ seconds the resampling operation touches $T f_s$ samples once, and the
channel reduction is a single pass over the same data. The wavelet
decomposition of Equation~\eqref{eq:mra} is likewise linear: a fast wavelet
transform to level $J$ performs filtering and downsampling at each level, and
because each level operates on half the samples of the one before, the total
work is bounded by
\begin{equation}
\sum_{j=0}^{J-1} c \cdot \frac{T f_s}{2^{j}} < 2 c \, T f_s,
\label{eq:cost-dwt}
\end{equation}
where $c$ is the cost per sample of the filter pair. The decomposition
therefore costs less than twice a single filtering pass regardless of $J$,
which is why increasing the decomposition level is inexpensive and why the
choice of $J$ was governed by the band structure of the signal rather than by
cost. Thresholding and reconstruction are of the same order.

Spectrogram extraction is dominated by the transforms. With a window of
$L$ samples advancing by $L/2$, a recording yields approximately
$2 T f_s / L$ frames, each requiring a transform of cost
$\mathcal{O}(L \log L)$, so the total is $\mathcal{O}(T f_s \log L)$. The
perceptual representation adds the filter bank, which applies $M$ triangular
weightings to each frame at a cost proportional to $M L / 2$ per frame, and
the subsequent resize to the network input size. These are small relative to
the transforms themselves.

Inference dominates everything preceding it. A forward pass through a
transferred architecture on a $224 \times 224$ input requires on the order of
$10^8$ to $10^9$ multiply accumulate operations depending on the model, which
exceeds the cost of the entire preprocessing chain for a recording of a few
seconds by two or more orders of magnitude. Two consequences follow. The
first is that the denoising stage should be judged on whether it improves the
result rather than on its cost, which is negligible in comparison. The second
is that the choice of architecture is the decision that determines whether the
system can run continuously, which is the reason a capacity constrained model
is included in the comparison at all and why its result is reported alongside
the most accurate one in Section~\ref{sec:res-deploy} rather than being
treated as an also ran.

\section{Experimental setup}\label{sec:setup}

This section gives the corpus and its partition, the training configuration,
the metrics, and the information required to repeat the experiments.
Section~\ref{sec:s-data} describes the data, Section~\ref{sec:s-train} the
optimisation, Section~\ref{sec:s-metrics} the measures reported, and
Section~\ref{sec:s-repro} the reproducibility details.

\subsection{Corpus and partition}\label{sec:s-data}

Experiments use COUGHVID, an open corpus of crowdsourced recordings collected
through a web interface~\cite{orlandic2021}. Contributors submitted audio from
their own devices in their own surroundings, with no control over microphone
quality, distance, background noise or recording level, and a substantial
fraction of submissions contain no cough. That heterogeneity is the reason the
corpus is appropriate for the present question: a detector evaluated on
studio material would not be tested against the conditions that make
denoising worth considering.

All recordings were converted to WAV and assigned to one of two classes,
cough or noncough, and partitioned into training and test sets.
Table~\ref{tab:split} gives the composition. The partition is fixed across all
eight configurations, so every model is trained on the same recordings and
measured on the same held out material, and differences between configurations
cannot be attributed to a favourable split.

\begin{table}[t]
\centering
\footnotesize
\caption{Composition of the partition, identical for all eight
configurations. The imbalance of roughly 2.15 to 1 in favour of cough is a
property of the collection rather than of the partition, and
Section~\ref{sec:res-errors} shows that it inflates accuracy relative to
measures that weight the classes equally.}
\label{tab:split}
\begin{tabular}{@{}lrrr@{}}
\toprule
Partition & Cough & Noncough & Total \\
\midrule
Training & 12{,}495 & 5{,}808 & 18{,}303 \\
Test & 3{,}124 & 1{,}453 & 4{,}577 \\
\midrule
Total & 15{,}619 & 7{,}261 & 22{,}880 \\
\bottomrule
\end{tabular}
\end{table}

The imbalance deserves comment because it conditions how the results in
Section~\ref{sec:results} should be read. Coughs outnumber noncoughs by
slightly more than two to one in both partitions, which reflects the fact that
contributors to the corpus were asked to cough. A classifier that labelled
every test recording as cough would therefore be correct on 3,124 of 4,577
occasions, an accuracy of 68.26 percent, without having learned anything at
all. Every accuracy figure reported below should be read against that
reference rather than against 50 percent.

\subsection{Training configuration}\label{sec:s-train}

All models were trained in MATLAB using the Deep Learning Toolbox on a
workstation with an NVIDIA RTX 4050 graphics processor. Optimisation used
Adam with a batch size of 32 throughout, a learning rate of $10^{-3}$ for
every model except DenseNet201, and $10^{-4}$ for DenseNet201.
Table~\ref{tab:train} records the settings, including the three points on
which the configurations are not uniform.

\begin{table}[t]
\centering
\footnotesize
\caption{Training configuration. The epoch budget and the learning rate are
the two respects in which the transferred models were not treated
identically, and the material held out for measurement was also the material
monitored during training. Section~\ref{sec:res-limits} states what each of
these costs the interpretation of the results.}
\label{tab:train}
\begin{tabular}{@{}ll@{}}
\toprule
Setting & Value \\
\midrule
Framework & MATLAB Deep Learning Toolbox \\
Hardware & NVIDIA RTX 4050 \\
Optimiser & Adam \\
Learning rate, all models except DenseNet201 & $10^{-3}$ \\
Learning rate, DenseNet201 & $10^{-4}$ \\
Batch size & 32 \\
Epochs, compact network, ResNet18, MobileNetV2 & 10 \\
Epochs, DenseNet201 & 3 \\
Monitoring interval & every 50 iterations (ResNet18, MobileNetV2) \\
 & every 200 iterations (DenseNet201) \\
Data used for monitoring & the test partition \\
Network retained & the one at the final iteration \\
Random seed & not fixed \\
Sampling rate, perceptual axis models & 16~kHz, mono \\
Sampling rate, compact network & as recorded, mono \\
Wavelet & sym4, level 6, empirical Bayes \\
Linear axis spectrogram & 256 point window, 200 overlap, \\
 & 256 point transform \\
Perceptual spectrogram & 512 point window, 256 overlap, 128 filters \\
Network input size & $224 \times 224$ \\
\bottomrule
\end{tabular}
\end{table}

DenseNet201 was trained for three epochs against ten for the other two
transferred models, because its depth makes each epoch considerably more
expensive, and at a learning rate an order of magnitude lower. These are real
asymmetries in the experimental design and we do not present them as anything
else. Their consequence is discussed in Section~\ref{sec:res-limits}, where we
set out which of our conclusions they can and cannot affect.

Two further properties of the procedure must be stated plainly, because they
bear on how every number in Section~\ref{sec:results} should be read. The
first is that no separate validation partition was constructed. The test
partition was supplied to the training routine as the data against which
progress was monitored, and the accuracy reported for each configuration is
measured on that same material. No quantity computed on the test partition
was used to select weights, since the network retained is the one at the final
iteration rather than the one at the best monitored score, but the
arrangement removes the separation between the data used to observe training
and the data used to report performance. The figures below are therefore not
estimates of performance on unseen material in the strict sense, and we treat
them as measurements on a partition that was visible throughout rather than as
generalisation estimates.

The second is that the compact network and the three transferred networks do
not receive comparable inputs, and the difference is larger than the change of
frequency axis alone. For the transferred models the perceptual spectrogram is
computed as a matrix, compressed logarithmically, mapped to the unit interval
and written as a single channel image, which is then replicated across three
channels at load time. For the compact network the linear axis spectrogram is
drawn to a figure and the rendered figure is captured as a raster image, so
the network receives a picture of a plot rather than the transform itself,
carrying the display colour map and whatever resolution the rendering surface
provided. The two paths also differ in sample rate, because the recordings are
resampled to 16~kHz only on the perceptual axis branch. The comparison between
the compact network and the transferred networks therefore varies five things
at once, and Section~\ref{sec:res-limits} explains why we decline to attribute
the difference between them to any one of the five.

\subsection{Evaluation metrics}\label{sec:s-metrics}

Let TP and TN denote correctly classified cough and noncough recordings, and
FP and FN the two error types, where a false positive is a noncough recording
labelled as a cough. The primary measures are
\begin{align}
\text{Accuracy} &= \frac{\mathrm{TP} + \mathrm{TN}}
{\mathrm{TP} + \mathrm{TN} + \mathrm{FP} + \mathrm{FN}}, \label{eq:acc} \\[4pt]
\text{Precision} &= \frac{\mathrm{TP}}{\mathrm{TP} + \mathrm{FP}},
\label{eq:prec} \\[4pt]
\text{Recall} &= \frac{\mathrm{TP}}{\mathrm{TP} + \mathrm{FN}},
\label{eq:rec} \\[4pt]
F_1 &= \frac{2 \cdot \text{Precision} \cdot \text{Recall}}
{\text{Precision} + \text{Recall}}. \label{eq:f1}
\end{align}

Precision and recall are reported alongside accuracy because the two error
types carry unequal cost in this application. A false negative is a cough that
goes unrecorded, which in a monitored period may be recovered when the person
coughs again. A false positive admits audio containing no cough into a
downstream analysis that was not trained on such material, and the resulting
output is reported with a confidence value it has not earned. Precision is
therefore the quantity a screening deployment watches most closely, and
Section~\ref{sec:res-errors} examines how it is composed for the leading
configuration.

Because the partition is imbalanced, accuracy alone would be a misleading
summary, and all four measures are reported for every configuration so that a
gain arising from a shift toward the majority class can be distinguished from
a genuine improvement. The relationship between the four rates and the
underlying counts is set out in Appendix~\ref{app:derivation}.

\subsection{Reproducibility}\label{sec:s-repro}

A single run consists of the following steps, performed end to end for one
configuration. Every recording in the corpus is read, reduced to mono and
resampled to 16~kHz. Where the configuration calls for it, the waveform is
denoised by a level six sym4 decomposition with soft thresholding and
reconstructed. The conditioned waveform is converted to the spectrogram
representation the configuration specifies and, for the transferred models,
resized to $224 \times 224$. The classifier is initialised, either randomly or
from transferred weights with a replaced two unit head, and optimised with the
settings of Table~\ref{tab:train} over the 18,303 training recordings. The
trained model is then applied once to the 4,577 test recordings and the
confusion counts recorded, from which the rates reported in
Table~\ref{tab:main} are computed.

The experimental grid is the full cross of two factors over four
architectures, with the constraint that the compact network is paired only
with the linear axis representation and the transferred models only with the
perceptual one. This yields eight runs: the compact network with and without
denoising, and each of ResNet18, MobileNetV2 and DenseNet201 with and without
denoising. Each run is trained once, so the figures reported carry no estimate
of variation across random initialisations, a limitation addressed in
Section~\ref{sec:res-limits}.

Two secondary analyses are reported alongside the primary results. The first
takes the difference between the paired runs of each architecture, answering
the question of whether denoising acts in the same direction and by the same
mechanism across the four models rather than merely whether it raises
accuracy; its result appears in Table~\ref{tab:delta}. The second examines
the confusion counts of the two ResNet18 configurations, answering the
question of how the improvement attributed to denoising is distributed
between the two error types; its result appears in Table~\ref{tab:resnet}.
Both use quantities recorded during the
evaluation runs described above and require no additional training. The
arithmetic relating the counts to the summary rates is given in
Appendix~\ref{app:derivation}, so that every figure in the body can be
checked against every other.
